\scanpage{30}
% archive sheet number: 29
et notre formule dans ce cas se réduit à
\[
\boxed{i(\overline{C'_{\eta}},\overline{D'_{\eta}})=
 \langle C'_{\eta},D'_{\eta}\rangle_v}
\]
\margin{\deleted{si $C',D'$ sans \unsure{composantes}}}

Or \inserted{dans} le cas général, on \unsure{voit|sait} aussitôt qu'on aura
\[
D_{\eta}\sim\alpha^*(D'_{\eta}),
\]
où $D'_{\eta}$ est un diviseur sur $A_{\eta}$ tel que
\[
{\overset{\deleted{-}}{D'_{\eta}}}\not\supset\alpha_{\eta}(X_{\eta})
\quad\text{et}\quad
{\overline{D'_{\eta}}}_0\not\supset\alpha_0(X_0)
\]
[Lemme : \unsure{dégager un|dégager en}]

d'où
\[
D=\alpha^*(D')+\operatorname{div}(f),
\]
où $D'$ est tel que $\alpha^*(D')$ soit défini et sans composante\deleted{\unsure{s}} verticale, et $f\in R^*(X)$. \deleted{\unsure{On|Donc}}

On est donc ramené à prouver la formule
\[
\begin{aligned}
&\boxed{\langle C_{\eta},\operatorname{div}f_{\eta}\rangle_v=i(C.\operatorname{div}f)}\\[-2pt]
&\qquad\rotatebox{90}{$=$}\\[-2pt]
&\quad v(f_{\eta}(C_{\eta}))
\end{aligned}
\]
\margin{valable si $\operatorname{div}(f)$ est sans \unsure{composantes|composante} \unsure{horizontales|horizontale}}
